%% ============================================================
%% CHAPTER 6 — CONCLUSION AND FUTURE WORK
%% ============================================================

\chapter{Conclusion and Future Work}
\label{ch:conclusion}

\section{Summary of Contributions}
\label{sec:conclusion_summary}

This thesis has developed, implemented, and validated the first
end-to-end physics-informed game-theoretic security framework for
gas transmission pipeline infrastructure.  The four original
contributions, summarised below, address a gap that spans three
decades of parallel but unconnected developments in structural
integrity, security game theory, and adversarial machine learning.

\paragraph{Contribution 1 — Physics-Informed Payoff Engine.}
The Monte Carlo failure probability engine, grounded in
\ac{BS7910}:2019 Level~2 \ac{FAD} and \ac{IIW} fatigue standards and
calibrated against 216 \ac{PHMSA} weld failure events, produces
$\Pf \in [\num{0.29}, \num{0.93}]$ across the test network.  The
resulting hierarchy (seamless $< $ DSAW $<$ HF-ERW $<$ LF-ERW $<$
fillet) reproduces the empirical PHMSA failure-rate ordering, validating
the calibration methodology.

\paragraph{Contribution 2 — Bayesian Stackelberg SSE.}
The LP-enumeration solver with three attacker types delivers a
physics-informed \ac{SSE} that outperforms uniform allocation by
\SI{17.0}{\percent} in total expected network loss at a
\SI{30}{\percent} defender budget.  The budget sensitivity analysis
demonstrates diminishing returns beyond $B/N \approx \num{0.60}$,
providing an evidence-based decision boundary for resource planning.

\paragraph{Contribution 3 — Adversarial NDE Threat Model.}
The \textsc{WeldDefectMLP} classifier (99.7\% clean accuracy) is shown
to be vulnerable to white-box $L_\infty$ attacks: \ac{BIM} achieves
\SI{20.1}{\percent} \acs{ASR} at $\eps = \num{0.30}$.  The
physics-informed $\eps$ scaling (\cref{eq:eps_scaling}) reveals that
the most game-theoretically prioritised segments face the highest
adversarial NDE threat, identifying a compounding risk unaddressed
by either standard alone.

\paragraph{Contribution 4 — Integrated Dashboard.}
The three-tab Plotly Dash application unifies all four model layers
into a deployable decision-support tool with click-to-inspect FAD
assessment, SSE coverage visualisation, and scenario comparison.
The 310-test suite with full public-interface coverage ensures
reproducibility and extensibility.

\section{Implications for Mediterranean Pipeline Security}
\label{sec:conclusion_med}

The STRATEGOS programme's focus on Mediterranean energy security makes
the extension to European corridors the most immediate application.
The Trans-Adriatic Pipeline (\ac{TAP}, \SI{878}{\kilo\meter},
\SI{1.2}{\meter} DSAW, \acs{ASME}~B31.8 design basis) and the
SNAM~Rete~Gas onshore network (Italy, \SI{32000}{\kilo\meter},
mixed vintage) share the physical characteristics of the synthetic
network: heterogeneous seam types, legacy ERW segments, and high
betweenness nodes at compressor stations.

The calibration pipeline is directly transferable: European regulators
(ENTSOG, ANSFR-T, ARERA) maintain incident databases analogous to
\ac{PHMSA}, and \ac{BS7910} is the mandated standard for FFS
assessment across EU member state operators.  The primary adaptation
required is re-calibrating the SCF distributions against European
incident data and extending the network topology to real GIS
infrastructure, which can be imported directly from ENTSOG's
publicly available TEN-E network data.

\section{Limitations}
\label{sec:conclusion_limitations}

\begin{enumerate}
  \item \textbf{Synthetic topology.}  The Gulf Coast network is
    representative but not operator-specific.  Conclusions about
    absolute risk levels should not be transferred to real systems
    without re-calibration.

  \item \textbf{Feature-vector NDE model.}  The adversarial attacks
    operate in a 32-dimensional feature space rather than directly
    on radiograph images.  While the threat model is conceptually
    valid, the ASR values are not directly comparable to attacks on
    real image-based NDE systems.

  \item \textbf{Single-stage game.}  The Bayesian SSE is a one-shot
    model.  Dynamic games with learning adversaries (e.g., repeated
    Stackelberg with Bayesian updating) are not addressed.

  \item \textbf{Correlated failures.}  The \ac{MC} engine treats
    segment failures as independent.  Correlated failure mechanisms
    (e.g., pressure surge, soil settlement) are outside scope.
\end{enumerate}

\section{Directions for Future Research}
\label{sec:conclusion_future}

Six directions for future research emerge directly from the
limitations above:

\begin{enumerate}
  \item \textbf{Image-based CNN adversarial NDE.}  Replace
    \textsc{WeldDefectMLP} with a small CNN (e.g., 3 conv layers)
    trained on synthetic 64~$\times$~64 weld radiograph images,
    implementing FGSM/PGD via PyTorch autograd for a higher-fidelity
    adversarial threat model.

  \item \textbf{Real network topologies.}  Import ENTSOG GeoJSON
    pipeline data (available at \url{https://www.entsog.eu}) to
    replace the synthetic Gulf Coast topology with a realistic
    European corridor, maintaining the same calibration pipeline.

  \item \textbf{Dynamic Stackelberg with adversary learning.}  Extend
    the SSE to a repeated game where the defender observes attack
    history and updates attacker type priors via Bayesian
    inference~\cite{Zhu2013}.

  \item \textbf{Correlated Monte Carlo.}  Incorporate spatial
    correlation in material properties (e.g., through a
    Gaussian random field on yield strength) to model batch
    manufacturing defects and soil-condition clustering.

  \item \textbf{Certified adversarial robustness.}  Apply interval
    bound propagation or randomised smoothing~\cite{Carlini2017}
    to provide provable robustness guarantees for the NDE classifier
    at specified $\eps$ levels, enabling standards-based acceptance
    criteria for AI-assisted inspection.

  \item \textbf{Cyber-physical attack modelling.}  Extend the
    adversarial layer to include SCADA signal manipulation (pressure
    and flow sensor spoofing), connecting the NDE threat model to
    the broader cyber-physical security literature~\cite{Mo2012}.
\end{enumerate}

\section{Closing Remarks}
\label{sec:conclusion_closing}

The central argument of this thesis is that pipeline security cannot
be reduced to either pure structural analysis or pure strategic modelling
independently: the physics determines what is worth attacking, and the
adversary's knowledge of the physics determines how they will attack.
The $\Pf$-informed Stackelberg game captures the first relationship;
the adversarial NDE threat model captures the second.  Together they
define a threat surface that is quantifiably larger than either
framework alone would suggest.

The \SI{17.0}{\percent} reduction in network expected loss achieved
by physics-informed SSE over uniform allocation, and the \SI{20.1}{\percent}
\ac{BIM} \acs{ASR} at $\eps = \num{0.30}$, are not large numbers in
isolation.  Their significance lies in what they make tractable:
the systematic, standards-grounded quantification of a cross-domain
threat that has so far been addressed only by expert judgment, and
the provision of a deployable decision-support tool for operators
and regulators who must make resource-allocation decisions under
precisely that uncertainty.
